% ファイル名は「番号_タイトル.tex」
% 番号はNotionを参考
\newpage
\section{共通点探し}  % 割り振りの名前を記入

\subsection{場所・人員配置}
\begin{table}[h]
  \centering
  \begin{minipage}[t]{0.48\columnwidth}
    \centering
    \caption{場所・該当グループ}
    % グループは2/10に分けたやつです．企画書班のNotionにも追加しておくので参考にしておいてください．
    \begin{tabular}{cl}
      \hline
      場所 & K101 \\
      \hline
      グループ & 全員 \\
      \hline
    \end{tabular}
  \end{minipage}
  \begin{minipage}[t]{0.48\columnwidth}
    \centering
    \caption{役割分担}
    \begin{tabular}{ll}
      \hline
      \multicolumn{1}{c}{役割} & \multicolumn{1}{c}{担当者} \\
      \hline
      司会 & 髙橋蒼 \\
      タイムキーパー & 佐藤謙成 \\
      スライド管理 & 細川夏風 \\
      班付き &  \\
      \hline
    \end{tabular}
  \end{minipage}
\end{table}

\subsection{詳細タイムスケジュール}
\vspace{0.5em}
\fontsize{8pt}{12pt}\selectfont

\begin{longtable}[c]{|p{.04\textwidth}|p{.2\textwidth}|p{.25\textwidth}|p{.45\textwidth}|}
  \caption{スケジュール} \\
  \hline
  \multicolumn{1}{|c|}{時間} & \multicolumn{1}{c|}{担当} & \multicolumn{1}{c|}{内容} & \multicolumn{1}{c|}{備考} \\
  \hline
  \endfirsthead
  
  \hline
  \multicolumn{1}{|c|}{時間} & \multicolumn{1}{c|}{担当} & \multicolumn{1}{c|}{内容} & \multicolumn{1}{c|}{備考} \\
  \hline
  \endhead
  
  \hline
  \endfoot
  
  \hline
  \endlastfoot
  
  \begin{tpacol}
    13:40
  \end{tpacol}&
  \begin{tpbcol}
    班付き
  \end{tpbcol}&
  \begin{tpccol}
    共通点探しに参加
  \end{tpccol}&
  \begin{tpdcol}
  \end{tpdcol}\\

  \begin{tpacol}
    
  \end{tpacol}&
  \begin{tpbcol}
    司会
  \end{tpbcol}&
  \begin{tpccol}
    イベントを回す
  \end{tpccol}&
  \begin{tpdcol}
  \end{tpdcol}\\

  \begin{tpacol}
    
  \end{tpacol}&
  \begin{tpbcol}
    タイムキーパー
  \end{tpbcol}&
  \begin{tpccol}
    時間管理する
  \end{tpccol}&
  \begin{tpdcol}
  \end{tpdcol} \\

  \begin{tpacol}
    
  \end{tpacol}&
  \begin{tpbcol}
    スライド管理
  \end{tpbcol}&
  \begin{tpccol}
    スライド管理する
  \end{tpccol}&
  \begin{tpdcol}
  \end{tpdcol} \\

\end{longtable}
\normalsize

\subsection{準備物品}
\begin{table}[h]
  \centering
  \caption{物品}
  \begin{tabular}{lrl}
    \hline
    \multicolumn{1}{c}{物品} & \multicolumn{1}{c}{数量} & \multicolumn{1}{c}{使用用途・備考} \\
    \hline
    マイク & 1 & 司会 \\
    \hline  
  \end{tabular}
\end{table}

\subsection{備考}
% 特に書くことなければ省略してOK